\chapterOrPart{About}\label{ch:about}
This appendix contains a very short description of \LaTeX\ and explains the content and use of this template.

\section{About \LaTeX}\label{sec:aboutlatex}
\LaTeX\ is software to make high-quality technical and scientific documents. It is not a "what you see is what you get" (WYSIWYG) editor such as MS Word or OpenOffice. It is a typesetting engine, which means that it works more like a programming language.

\LaTeX\ requires an input file that contains the source code. The structure of that file is shown below and is discussed in more detail \href{https://en.wikibooks.org/wiki/LaTeX/Document_Structure}{here}.

\begin{verbatim}
    \documentclass[options]{class}
        This part is called the PREAMBLE, here:
        - packages are loaded that enhance the capabilities of LaTeX
        - settings are defined such as: margins, fonts, line spacing, ...
        
    \begin{document}
        This part is called the DOCUMENT, here:
        - content is provided in plain text format
    \end{document}
\end{verbatim}

\subsection{Document class}
The input file starts with the command \textbackslash\verb|documentclass[options]{class}|, where the parameter \texttt{class} stands for the name of a certain class file, recognised by the extension \texttt{.cls}. In a class file the layout standard is defined which tells \LaTeX\ how to format your content. The parameter \texttt{options} customizes the behaviour of the document class by setting parameters such as font size, paper size, etc.

This template uses the class file \texttt{memoir} but many other class files exist such as \texttt{article}, \texttt{report}, \texttt{book},\dots. A comprehensive list of class files can be consulted \href{https://ctan.org/topic/class}{at CTAN}.

\subsection{Preamble}
The next part in the input file is called the \texttt{preamble}. In this part, packages are loaded that enhance the capabilities of \LaTeX. This is typically done in so-called style files with the extension \texttt{.sty}. The packages that are loaded in this template can be found in the style file \texttt{styles/myPackages} and \texttt{styles/myListings}. The use and purpose of some of the loaded packages are discussed further in this chapter. Currently there exist around 4000 packages for \LaTeX\ for all kind of applications. The comprehensive list of all packages can be consulted \href{https://www.ctan.org/pkg/}{at CTAN}.

In the \texttt{preamble} a user can also define his/her own commands and environments. In this template, this is done in the other style files in the folder \texttt{styles}.

\subsection{Entering content}
The actual content of your work is enclosed between the commands \textbackslash\verb|begin{document}| and \textbackslash\verb|end{document}|. When you have a lot of content, you will typically divide it over several files. Files with content have the extension \texttt{.tex}. Content is added as plain text, combined with embedded \LaTeX-commands for specific typeset results. To include files in your input file, you can use the command: \textbackslash\verb|include{filename}| or \textbackslash\verb|input{filename}|.

The command \textbackslash\verb|include{filename}| is used to include chapters. It will process all the floats (figures and tables) from the previous chapter and cause a page break so that the next chapter can start on a new page. It will also start a new auxiliary file in the background to keep track of the locations of figures, tables and equation within the chapter for proper cross-referencing. Inspect the file \texttt{mainReport.tex} to see what chapters are included in this template. 

The command \textbackslash\verb|input{filename}| only imports the content from filename.tex into the target file. It's equivalent to typing all the content from filename.tex into the current file at the location of the \textbackslash\verb|input{filename}| command.

The content is structured for the reader by using commands such as \textbackslash\verb|chapter{name}|,\\ \textbackslash\verb|section{name}|, \textbackslash\verb|subsection{name}|, etc.

More info on how to add content such as text, images, tables and math can be found in \href{https://www.overleaf.com/learn/latex/Learn_LaTeX_in_30_minutes}{Learn \LaTeX\ in 30~minutes}.

\subsection{Compiling your input file}
You have to \texttt{compile} you input file to obtain a typeset PDF document. The processing of your input file is done by a so-called TeX engine. It will parse your input file, line by line, detect the embedded \LaTeX-commands in your input file and act upon them, according the settings  in your class file and packages.

This means that you only have to concentrate on the content, while /LaTeX/ will take care of the formatting!


\section{About this template}
\subsection{Content of the template}
This template will get you started with writing your thesis, report or slides in \href{https://www.latex-project.org}{\LaTeX}. The template contains:
\begin{itemshort}
    \item a title page for VUB and for BRUFACE
    \item acknowledgements
    \item an abstract in English, Dutch and French
    \item a table of contents
    \item a list of figures, tables and listings
    \item a nomenclature
    \item the typical chapters
    \item appendices with useful information on various topics
    \item a bibliography
\end{itemshort}

\subsection{Content of the appendices}
In the appendices you will find useful information about
\begin{itemshort}
    \item \LaTeX\ and this template
    \item formatting of figures, tables, equations and source code
    \item useful \LaTeX\ packages and commands
    \item examples of different types of figures and plots
    \item mathematical and Greek symbols
\end{itemshort}

The template uses the class \href{https://ctan.org/pkg/memoir?lang=en}{memoir}. It is a class for typesetting mathematical works as books, reports and articles.

\subsection{How to start using it}
This template is read-only: this means that you must copy it, before you start using it. A \href{https://vub.cloud.panopto.eu/Panopto/Pages/Sessions/List.aspx?folderID=eb5cbed1-a7f8-4d08-bca3-ad1e00db8d7a}{set of videos} is available that demonstrate the use of this template\footnote{The videos can only be accessed from a VUB-account. Some videos are Dutch spoken, others are English spoken}.

This template can generate a report and slides from the same source code. That's why there are two main files:
\begin{itemize}
    \item \texttt{mainReport.tex}: to compile a report or thesis
    \item \texttt{mainSlides.tex}: to compile the slides
\end{itemize}

Only selected content is put on the slides though. To see all content in this template, you must compile \texttt{mainReport.tex} and/or inspect its source code.

\subsection{How to change the language}

This template is written in English. If you write in another language, you have to adapt the language in three ways (example below is for Dutch):
\begin{enumerate}
    \item Select \texttt{Dutch} as spell check language in Overleaf, via the \texttt{menu} in the upper left corner. Words that do not belong to the selected spell check language will be underlined in red in the editor. 
    \item Select the option \texttt{dutch} in \texttt{styles/myPackages.sty} in these two packages:
    \begin{itemize}
        \item \textbackslash\verb|usepackage[dutch]{babel}|\\
        This will automatically translate standard names:\\
        chapter $\to$ hoofdstuk\\
        table of content $\to$ inhoudsopgave\\
        bibliography $\to$ bibliografie\\\dots
        \item \textbackslash\verb|usepackage[dutch]{cleveref}|\\
        This will automatically translate cross-referencing names:
        \\figure $\to$ figuur\\
        table $\to$ tabel\\
        equation $\to$ vergelijking\\\dots
    \end{itemize}
\end{enumerate}
    

\mode<all>